% Part: many-valued-logic
% Chapter: sequent-calculus
% Section: rules-and-proofs

\documentclass[../../../include/open-logic-section]{subfiles}

\begin{document}

\olfileid[fr]{mvl}{seq}{rul}

\olsection{Règles et \usetoken{p}{derivation}}

Dans un séquent, $\Gamma, \Delta, \Pi, \Lambda$ désignent des suites
finies de !!{formula}s. Dans une logique à $n$ valeurs, nous fixons un
ordre sur les valeurs de vérité. La place~$i$ correspond à la
$i$-ième valeur.

\begin{defn}[Séquent]
Un \emph{séquent à $n$ places} est une expression de la forme
\[
\Gamma_1 \nSequent \dots \nSequent \Gamma_n
\]
où chaque $\Gamma_i$ est une suite finie, éventuellement vide,
de !!{formula}s du langage~$\Lang L$.
\end{defn}

\begin{defn}[Séquent initial]
Un \emph{séquent initial à $n$ places} est un séquent à $n$ places de la
forme $!A \nSequent \dots \nSequent !A$, pour toute !!{formula}~$!A$
du langage.

Si le langage contient un connecteur d'arité~$0$, c'est-à-dire une
constante propositionnelle~$\star$, nous admettons aussi comme séquent
initial $\dots \nSequent \star \nSequent \dots$ : $\star$ figure
uniquement à la place correspondant à sa valeur de vérité
$\tf{\star}\in V$, les autres places étant vides.
\end{defn}

Pour chaque connecteur d'une logique à $n$ valeurs~$\Log{L}$, il y a
une règle logique pour chacune des valeurs que ce connecteur peut
prendre dans~$\Log{L}$. Les !!{derivation}s d'un calcul des séquents
à $n$ places pour~$\Log{L}$ sont des arbres de séquents dont les
séquents situés tout en haut sont initiaux. Tout séquent placé sous un
ou plusieurs autres doit s'en déduire correctement par une règle
d'inférence relative aux connecteurs de~$\Log{L}$.

\begin{defn}[Théorèmes]
!!^a{formula}~$!A$ est un \emph{théorème} d'une logique à $n$
valeurs~$\Log{L}$ s'il existe !!a{derivation} du séquent à $n$ places
qui contient~$!A$ à chacune des places correspondant à une valeur
désignée de~$\Log{L}$ et dont les autres places sont vides. Nous
écrivons $\Proves[\Log{L}] !A$ si $!A$ est un théorème, et
$\Proves/[\Log{L}] !A$ dans le cas contraire.
\end{defn}

\begin{defn}[!!^{derivability}]
!!^a{formula}~$!A$ est \emph{!!{derivable} de} l'ensemble
de !!{formula}s~$\Gamma$ dans une logique à $n$ valeurs~$\Log{L}$,
ce que nous notons $\Gamma\Proves[\Log{L}] !A$, s'il existe un
sous-ensemble fini~$\Gamma_0\subseteq\Gamma$ et une suite~$\Gamma_0'$
énumérant les !!{formula}s de~$\Gamma_0$ tels que le séquent suivant
admette !!a{derivation} :
\[
\Lambda_1 \nSequent \dots \nSequent \Lambda_n.
\]
Ici, $\Lambda_i$ est $!A$ si la place~$i$ correspond à une valeur
désignée; dans le cas contraire, $\Lambda_i$ est~$\Gamma_0'$.
Si $!A$ n'est pas !!{derivable} de~$\Gamma$, nous écrivons
$\Gamma\Proves/[\Log{L}] !A$.
\end{defn}

Par exemple, la logique à trois valeurs de \L ukasiewicz possède un
calcul des séquents à trois places. Dans un séquent
$\Gamma \nSequent \Pi \nSequent \Delta$, $\Gamma$ correspond
à~$\False$, $\Pi$ à~$\Undef$ et $\Delta$ à~$\True$. Les axiomes sont
$!A \nSequent !A \nSequent !A$. Puisque seule la valeur~$\True$ est désignée,
$\Gamma\Proves[\LogLuk[3]] !A$ si et seulement s'il existe une
suite~$\Gamma_0'$ énumérant un sous-ensemble fini de~$\Gamma$
telle que $\Gamma_0' \nSequent \Gamma_0' \nSequent !A$ admette
!!a{derivation}. Si $\Undef$ était aussi désignée, le séquent à
dériver serait $\Gamma_0' \nSequent !A \nSequent !A$, avec la même
condition sur~$\Gamma_0'$.

\end{document}
